\documentclass[a4paper,12pt,polish]{book}

\textwidth=161mm
\textheight=247mm
\evensidemargin=0mm 
\oddsidemargin=0mm
\topmargin=-12mm

%\documentclass[a4paper,12pt]{book}
\usepackage[utf8]{inputenc}
\usepackage[polish]{babel}
\usepackage[T1]{fontenc}
\usepackage{polski}
\usepackage{graphicx}
\usepackage{amsmath}
\usepackage{comment}
\usepackage{amsfonts}
\usepackage[font={small}]{caption}
\usepackage{float}
\usepackage{rotating}
\usepackage{color}

\usepackage{rotating}
%\usepackage{jneurosci}
\usepackage{setspace}
\usepackage{subfigure}
\newcommand{\g}[1]{\mathbf{#1}}

\usepackage{natbib}
%\usepackage{geometry}
%\geometry{bindingoffset=1cm}
%\evensidemargin=0mm 
%\oddsidemargin=0mm
\begin{document}
\onehalfspacing
%\doublespacing
\thispagestyle{empty}

\widowpenalty=1000 
\clubpenalty=500
 
\begin{center}
  
   \ \\
 

\includegraphics{../ryciny/nencki.png}
\vfill

   \begin{doublespace}
     \hspace*{-4ex}\begin{minipage}[l]{1.1\textwidth}
       \begin{center}
         {\sc\LARGE  
	Walidacja metod analizy rejestracji elektrofizjologicznych \\
 	przy użyciu danych symulacyjnych
      }\\[1em]
         {\large\sc Helena Głąbska}
       \end{center}

     \end{minipage}
   \end{doublespace}

   \vfill
\
 \vfill
  {%\hfill
  \begin{minipage}[c]{0.75\textwidth}
    \centering
    {\sc praca doktorska wykonana pod kierunkiem} \\
    {\sc dr. hab.\ Daniela Wójcika}\\
  \end{minipage}} \\
  \vfill  

  {\sc Instytut Biologii Doświadczalnej} \\
  {\sc im. Marcelego Nenckiego}\\
  {\sc Polska Akademia Nauk}\\[1em]
  {\sc Warszawa $\ast$ 2015}
\end{center}

\newpage
\thispagestyle{empty}
\ 
\newpage
\thispagestyle{empty}
\vspace*{0.5\textheight}

\hspace*{0.5\textwidth} \begin{minipage}[c]{0.5\textwidth}
  %\centering
  {\sc Podziękowania ...}
\end{minipage}

\ 
\newpage


\ 
\newpage





\section*{Streszczenie} 

\tableofcontents
\chapter{Wstęp}
\section{Modelowania układu nerwowego - modele typu Hodgkina-Huxleya}
\subsection{Model Hodgina-Huxleya}
\subsection{Rozszerzenie modelu Hodgina-Huxleya}
\subsection{Symulatory układu nerwowego}
\subsubsection{Neuron}
\subsubsection{GENESIS}
\subsubsection{NEST}
\subsubsection{Brian}
\subsubsection{STEPS}
\subsubsection{XPPAUT}
\section{Potencjał zewnątrzkomórkowy}
\section{Laminar Population Analysis (LPA)}
\section{Laminar Population Analysis (LPA)}
Laminar Population Analysis (LPA) \cite{Einevoll2007} jest jedną z metod analizy danych elektrofizjologicznych. Podobnie jak wiele innych popularnych metod, rozkłada oryginalny sygnał na składowe, które próbuje zinterpretować w kontekście aktywności poszczególnych populacji komórek. O ile w szeroko stosowanej metodzie PCA (Prinipal Compomnents Analysis) sygnał jest rozkładany na składowe główne, w ICA (Independent Components Analysis) na składowe niezależne, celem LPA jest wyodrębnienie wkładów wywołanych aktywnością poszczególnych populacji. Metoda LPA wychodzi z następujących założeń:
\begin{itemize}
\item sygnał LFP jest wywołany przez iglice potencjał czynnościowego, który możemy zaobserwować w sygnale MUA
\item LFP i MUA mają postać produktową,  (lub są sumą kilku składowych o postaci produktowej w gLPA)
\end{itemize}
\subsection{Moldelowanie MUA}
Pierwszy etap metody LPA polega na oszacowaniu tempa wystrzeliwania iglic potencjału czynnościowego. Sygnał MUA jest rozpisany na wkłady pochodzące od poszczególnych populacji:
\begin{equation}
MUA^m(z_i, t_j) = \sum^{N_{pop}}_{n=1} M_n(z_i)r_n(t_j).
\label{MUAdecomposition}
\end{equation}
W tym przypadku $N_{pop}$ to liczba populacji, $r_n(t_j)$ - tempo wystrzeliwania iglic potencjału czynnościowego,  $M_n$ - profil przestrzenny MUA dla n-tej populacji. $M_n$ zależy od własności fizjologicznych neuronów w danej populacji, ich rozmieszczenia oraz przewodności elektrycznej ośrodka. Metoda LPA przyjmuje że $M_n$ jest funkcją trapezoidalną, definiowaną za pomocą trzech parametrów: $z_{0n}, \ a_n, \ b_n$. Przyjmuje wartość 0 dla $z_{j}< z_{0n} - a_n/2 - b$ oraz $z_{j} > z_{0n} + a_n/2 + b$, wartość 1 dla $z_{0n} - a_n/2 < z_j < z_{0n} + a_n/2$, jest liniowa w pozostałych przedziałach (patrz fig. \ref{MUAshape}). Dodatkowo zakłada się że $b_n < 0.1 mm$.

W tym etapie wyznaczmy parametry $z_{0n}, \ a_n, \ b_n$, a tym samym determinujemy tempo wystrzeliwania iglic potencjału czynnościowego dla poszczególnych populacji $r_n$ minimalizując błąd \ref{MUAerror}.

\begin{figure}
\caption{MUA shape profile function}
\label{MUAshape}
\includegraphics{../ryciny/MUAshape.png}
\end{figure}

\begin{equation}
e_M = \sum^{N_{ch}}_{i=1}\sum^{N_t}_{j=1} \left( MUA(z_i, t_j) - MUA^m(z_i, t_j) \right)^2/ \sum^{N_{ch}}_{i=1}\sum^{N_t}_{j=1} MUA(z_i, t_j)^2.
\label{MUAerror}
\end{equation}

\subsection{Modelowanie LFP}
W tym kroku korzystamy z założenia że sygnał LFP jest wywołany przez iglice potencjału czynnościowego, który możemy zaobserwować w sygnale MUA oraz że LFP ma postać produktową. Sygnał LFP rozpisuje się jako:
\begin{equation}
 LFP^m(z_i, t_j) = \sum^{N_{pop}}_{n=1}L_n(z_i)(h*r_n)(t_j),
 \label{LFPdecomposition}
\end{equation}
gdzie symbol $*$ oznacza splot co zapisuje się jako:
\begin{equation}
(h_n*r_n)(t_j) = \sum_{k = -\infty}^{\infty}h(t_k)r_n(t_j - t_k)
\label{splot}
\end{equation}

$L_n$ to profil przestrzenny LFP wywołanego aktywnością populacji $n,$ $h_n$ to jądro która reprezentuje opóźnienie w wystąpieniu i propagacji sygnału LFP wywołanego aktywnością populacji $n:$


\begin{eqnarray}
\label{h_function}
h(t_j) = \frac{1}{\tau}e^{-(t_j - \Delta)/\tau}  &for \ t_j \geq 0\\
h(t_j) = 0 &for \ t_j <0.
\end{eqnarray}

Wyznaczamy $\tau,$ $\Delta,$ $L_n$ minimalizując:
\begin{equation}
e_L = \sum^{N_{ch}}_{i=1}\sum^{N_t}_{j=1} \left( LFP(z_i, t_j) - LFP^m(z_i, t_j) \right)^2/ \sum^{N_{ch}}_{i=1}\sum^{N_t}_{j=1} LFP(z_i, t_j)^2.
\label{LFPerror}
\end{equation}
\section{Rozszerzenie metody LPA -generalized Laminar Population Analysis (gLPA)}
\section{Cel pracy}
\subsection{Walidacja metody Laminar Population Analysis}
\subsection{??  Walidacje metody Independent Component Analysis}
\subsection{Publikacja bogatego zestawu danych na potrzeby walidacji innych metod analizy danych}


\chapter{Dane}
\section{Model pętli korowo-wzgórzowej Traub at al}
\section{Rozbudowa modelu Traub at al}
\section{Wygenerowane zestawy danych}
\subsection{Stymulacja prądem oscylacyjnym}
\subsection{Odpowiedź sieci na wejście z pojedynczej populacji}
\subsection{Blokada szybkich kanałów jonowych}


\chapter{Wyniki}
\section{Walidacja metody generalized Laminar Population Analysis}
\subsection{Uproszczony przypadek - zastąpienie pierwszej fazy gLPA danymi z PSTH}
\subsection{Oryginala metoda gLPA}
\subsection{Modyfikacja założeń dotyczących profilu przestrzennego MUA}


\chapter{Dyskusja}

\chapter{Podsumowanie wyników i wnioski}




\bibliographystyle{agsm}
\bibliography{doktorat}

\end{document}
